% harmlessness_FR.tex — WAC (Windows Artefacts Collector) harmlessness file, French edition.
% The English edition is harmlessness_EN.tex, in this folder: both are maintained together,
% section for section.
%
% Build (LuaLaTeX, as the editor's recipe; build artefacts removed after):
%     ../build-docs.sh harmlessness_FR.tex
% The cover page is ../cover_harmlessness_FR.png, derived from ../cover_FR.png by
% ../retitle-cover.py. The layout is ../wac-style.tex, shared with every WAC document.
%
% This file says WHAT WAC does to the examined machine, every trace included, and
% HOW that is verified. It is to be revised with every change that touches a
% system call, the build's guards or the harness's harmlessness check.

\documentclass[11pt,a4paper,oneside]{report}

% ---------------------------------------------------------------- language and labels
\usepackage[french]{babel}
\newcommand{\wacdoctitle}{Dossier d'innocuité}
\newcommand{\wacpageof}{sur}
\newcommand{\wacnotetitle}{À retenir}
\newcommand{\wacwarningtitle}{Limite}
% Layout shared with every WAC document.
% wac-style.tex — layout shared by every WAC document (user manuals, testing
% guides), French and English editions alike: one copy, so that they cannot
% drift apart.
%
% A document sets, BEFORE % wac-style.tex — layout shared by every WAC document (user manuals, testing
% guides), French and English editions alike: one copy, so that they cannot
% drift apart.
%
% A document sets, BEFORE % wac-style.tex — layout shared by every WAC document (user manuals, testing
% guides), French and English editions alike: one copy, so that they cannot
% drift apart.
%
% A document sets, BEFORE \input{../wac-style.tex}:
%   \usepackage[<language>]{babel}
%   \newcommand{\wacdoctitle}{...}      running header: "WAC — <title>"
%   \newcommand{\wacpageof}{...}        footer: "Page N <of> M"
%   \newcommand{\wacnotetitle}{...}     default title of the "note" box
%   \newcommand{\wacwarningtitle}{...}  default title of the "attention" box
% and AFTER it its \hypersetup (PDF title, author, subject).

% ---------------------------------------------------------------- fonts, language
\usepackage{fontspec}
\setmainfont{Noto Sans}
\setmonofont{DejaVu Sans Mono}[Scale=0.84]

% ---------------------------------------------------------------- layout
\usepackage[top=26mm,bottom=26mm,left=24mm,right=22mm,headheight=14pt]{geometry}
\usepackage{graphicx}
\graphicspath{{../}}
\usepackage{xcolor}
\usepackage{tikz}
\usepackage{booktabs,tabularx,longtable,array}
\usepackage{enumitem}
\usepackage{fancyhdr}
\usepackage{titlesec}
\usepackage[most]{tcolorbox}
\usepackage{listings}
\usepackage{url}
\usepackage{lastpage}
\usepackage[hidelinks,unicode,pdfpagelabels]{hyperref}

% Colours sampled from the cover page
\definecolor{marine}{HTML}{0A1A33}
\definecolor{bleu}{HTML}{2F6DB5}
\definecolor{accent}{HTML}{1F66D1}
\definecolor{fond}{HTML}{EEF3FA}
\definecolor{alerte}{HTML}{B3261E}
\definecolor{fondalerte}{HTML}{FBEDEC}
\definecolor{gris}{HTML}{5B6573}

\setlength{\parindent}{0pt}
\setlength{\parskip}{0.55em}
\linespread{1.08}
\setlist{itemsep=0.25em,topsep=0.35em}
\urlstyle{tt}
\renewcommand{\arraystretch}{1.25}
% Left-aligned table columns: justified narrow columns opened wide gaps
% between words.
\newcolumntype{L}[1]{>{\raggedright\arraybackslash}p{#1}}
% Monospaced file and field names cannot be hyphenated: let the line's spaces
% stretch instead.
\setlength{\emergencystretch}{3em}

% Headings
\titleformat{\chapter}[display]
  {\color{marine}\bfseries\raggedright}
  {\color{bleu}\large\MakeUppercase{\chaptertitlename}~\thechapter}
  {0.3em}
  {\Huge}[{\color{accent}\titlerule[2pt]}]
\titlespacing*{\chapter}{0pt}{0pt}{1.6em}
\titleformat{\section}{\color{marine}\Large\bfseries}{\color{bleu}\thesection}{0.8em}{}
\titleformat{\subsection}{\color{marine}\large\bfseries}{\color{bleu}\thesubsection}{0.7em}{}
\titleformat{\subsubsection}{\color{bleu}\normalsize\bfseries}{}{0em}{}

% Headers and footers
\pagestyle{fancy}
\fancyhf{}
\fancyhead[L]{\small\color{gris}WAC — \wacdoctitle}
\fancyhead[R]{\small\color{gris}\leftmark}
% Footer: "Page N of M" (the word set by \wacpageof), readable. A bare, small, light-grey number went
% unnoticed.
\newcommand{\piedpage}{\small\color{marine}\textbf{Page \thepage}\ \wacpageof\ \pageref*{LastPage}}
\fancyfoot[C]{\piedpage}
\renewcommand{\footrulewidth}{0.4pt}
\renewcommand{\footrule}{\hbox to\headwidth{\color{fond}\leaders\hrule height \footrulewidth\hfill}}
\renewcommand{\headrulewidth}{0.4pt}
\renewcommand{\headrule}{\hbox to\headwidth{\color{fond}\leaders\hrule height \headrulewidth\hfill}}
\renewcommand{\chaptermark}[1]{\markboth{#1}{}}
\fancypagestyle{plain}{\fancyhf{}\fancyfoot[C]{\piedpage}\renewcommand{\headrulewidth}{0pt}}

% Boxes
\newtcolorbox{note}[1][\wacnotetitle]{enhanced,breakable,colback=fond,colframe=bleu,
  boxrule=0pt,leftrule=3pt,arc=0pt,left=8pt,right=8pt,top=6pt,bottom=6pt,
  title={\bfseries #1},coltitle=bleu,colbacktitle=fond,attach title to upper={\par\smallskip}}
\newtcolorbox{attention}[1][\wacwarningtitle]{enhanced,breakable,colback=fondalerte,colframe=alerte,
  boxrule=0pt,leftrule=3pt,arc=0pt,left=8pt,right=8pt,top=6pt,bottom=6pt,
  title={\bfseries #1},coltitle=alerte,colbacktitle=fondalerte,attach title to upper={\par\smallskip}}

% Code
\lstset{basicstyle=\ttfamily\small,columns=fullflexible,keepspaces=true,
  backgroundcolor=\color{fond},frame=l,framerule=3pt,rulecolor=\color{bleu},
  xleftmargin=9pt,framexleftmargin=6pt,aboveskip=0.8em,belowskip=0.8em,
  breaklines=true,breakatwhitespace=false,showstringspaces=false,
  literate={é}{{\'e}}1 {è}{{\`e}}1 {à}{{\`a}}1 {ê}{{\^e}}1 {ô}{{\^o}}1 {ç}{{\c c}}1
           {É}{{\'E}}1 {—}{{---}}1 {’}{{'}}1 {…}{{...}}1 {«}{{\guillemotleft}}1 {»}{{\guillemotright}}1}

% Breakable monospace. A path or a field name in \texttt is one word that no
% hyphenation pattern cuts: in a table cell it ran 190 pt past the column
% (ProgramData\Microsoft\NetFramework\BreadcrumbStore). As the \seqsplit of
% the EventsManager theme's tables, a line may break inside it, without a
% hyphen that would read as part of the name: after \ / . _ - : and between a
% lowercase letter or digit and a capital (AmcacheSha1|Covers|Bytes). Only
% inside \texttt — the classes are set once, the breaks switched on there.
\newXeTeXintercharclass\wacBreakAfter
\newXeTeXintercharclass\wacUpper
\newXeTeXintercharclass\wacLower
\XeTeXcharclass`\\=\wacBreakAfter
\XeTeXcharclass`\/=\wacBreakAfter
\XeTeXcharclass`\.=\wacBreakAfter
\XeTeXcharclass`\_=\wacBreakAfter
\XeTeXcharclass`\-=\wacBreakAfter
\XeTeXcharclass`\:=\wacBreakAfter
\count255=`A \loop\XeTeXcharclass\count255=\wacUpper\ifnum\count255<`Z\advance\count255 by 1\repeat
\count255=`a \loop\XeTeXcharclass\count255=\wacLower\ifnum\count255<`z\advance\count255 by 1\repeat
\count255=`0 \loop\XeTeXcharclass\count255=\wacLower\ifnum\count255<`9\advance\count255 by 1\repeat
\XeTeXinterchartoks\wacBreakAfter\wacUpper={\allowbreak}
\XeTeXinterchartoks\wacBreakAfter\wacLower={\allowbreak}
\XeTeXinterchartoks\wacBreakAfter 0={\allowbreak}
\XeTeXinterchartoks\wacLower\wacUpper={\allowbreak}
% \textbackslash is no character token: it carries its own break.
\newcommand{\wacbackslash}{\char92\allowbreak}
\let\wacplaintexttt\texttt
\renewcommand{\texttt}[1]{\wacplaintexttt{\XeTeXinterchartokenstate=1\let\textbackslash\wacbackslash#1}}

% Shortcuts
\newcommand{\wac}{\textsc{wac}}
\newcommand{\opt}[1]{\texttt{-{}-#1}}
\newcommand{\champ}[1]{\texttt{#1}}
\newcommand{\fichier}[1]{\texttt{#1}}

% ============================================================ COVER PAGE
% The cover images are 2:3 (1024 x 1536), taller than an A4 page once widened
% to it: centred, they lost 9 mm at the top — the CEC logo — and 9 mm at the
% bottom. Raised 4 mm above the top of the page instead: the logo, which starts
% 6 mm below the image's edge, stays whole, and the bottom row of text keeps
% 7 mm to the page's edge; only a decorative corner below it is cut.
% Usage: \coverpage{cover_FR.png}
\newcommand{\coverpage}[1]{%
  \begin{titlepage}
  \thispagestyle{empty}
  \begin{tikzpicture}[remember picture,overlay]
    \node[inner sep=0pt,anchor=north,yshift=4mm] at (current page.north)
      {\includegraphics[width=\paperwidth]{#1}};
  \end{tikzpicture}
  \end{titlepage}}
:
%   \usepackage[<language>]{babel}
%   \newcommand{\wacdoctitle}{...}      running header: "WAC — <title>"
%   \newcommand{\wacpageof}{...}        footer: "Page N <of> M"
%   \newcommand{\wacnotetitle}{...}     default title of the "note" box
%   \newcommand{\wacwarningtitle}{...}  default title of the "attention" box
% and AFTER it its \hypersetup (PDF title, author, subject).

% ---------------------------------------------------------------- fonts, language
\usepackage{fontspec}
\setmainfont{Noto Sans}
\setmonofont{DejaVu Sans Mono}[Scale=0.84]

% ---------------------------------------------------------------- layout
\usepackage[top=26mm,bottom=26mm,left=24mm,right=22mm,headheight=14pt]{geometry}
\usepackage{graphicx}
\graphicspath{{../}}
\usepackage{xcolor}
\usepackage{tikz}
\usepackage{booktabs,tabularx,longtable,array}
\usepackage{enumitem}
\usepackage{fancyhdr}
\usepackage{titlesec}
\usepackage[most]{tcolorbox}
\usepackage{listings}
\usepackage{url}
\usepackage{lastpage}
\usepackage[hidelinks,unicode,pdfpagelabels]{hyperref}

% Colours sampled from the cover page
\definecolor{marine}{HTML}{0A1A33}
\definecolor{bleu}{HTML}{2F6DB5}
\definecolor{accent}{HTML}{1F66D1}
\definecolor{fond}{HTML}{EEF3FA}
\definecolor{alerte}{HTML}{B3261E}
\definecolor{fondalerte}{HTML}{FBEDEC}
\definecolor{gris}{HTML}{5B6573}

\setlength{\parindent}{0pt}
\setlength{\parskip}{0.55em}
\linespread{1.08}
\setlist{itemsep=0.25em,topsep=0.35em}
\urlstyle{tt}
\renewcommand{\arraystretch}{1.25}
% Left-aligned table columns: justified narrow columns opened wide gaps
% between words.
\newcolumntype{L}[1]{>{\raggedright\arraybackslash}p{#1}}
% Monospaced file and field names cannot be hyphenated: let the line's spaces
% stretch instead.
\setlength{\emergencystretch}{3em}

% Headings
\titleformat{\chapter}[display]
  {\color{marine}\bfseries\raggedright}
  {\color{bleu}\large\MakeUppercase{\chaptertitlename}~\thechapter}
  {0.3em}
  {\Huge}[{\color{accent}\titlerule[2pt]}]
\titlespacing*{\chapter}{0pt}{0pt}{1.6em}
\titleformat{\section}{\color{marine}\Large\bfseries}{\color{bleu}\thesection}{0.8em}{}
\titleformat{\subsection}{\color{marine}\large\bfseries}{\color{bleu}\thesubsection}{0.7em}{}
\titleformat{\subsubsection}{\color{bleu}\normalsize\bfseries}{}{0em}{}

% Headers and footers
\pagestyle{fancy}
\fancyhf{}
\fancyhead[L]{\small\color{gris}WAC — \wacdoctitle}
\fancyhead[R]{\small\color{gris}\leftmark}
% Footer: "Page N of M" (the word set by \wacpageof), readable. A bare, small, light-grey number went
% unnoticed.
\newcommand{\piedpage}{\small\color{marine}\textbf{Page \thepage}\ \wacpageof\ \pageref*{LastPage}}
\fancyfoot[C]{\piedpage}
\renewcommand{\footrulewidth}{0.4pt}
\renewcommand{\footrule}{\hbox to\headwidth{\color{fond}\leaders\hrule height \footrulewidth\hfill}}
\renewcommand{\headrulewidth}{0.4pt}
\renewcommand{\headrule}{\hbox to\headwidth{\color{fond}\leaders\hrule height \headrulewidth\hfill}}
\renewcommand{\chaptermark}[1]{\markboth{#1}{}}
\fancypagestyle{plain}{\fancyhf{}\fancyfoot[C]{\piedpage}\renewcommand{\headrulewidth}{0pt}}

% Boxes
\newtcolorbox{note}[1][\wacnotetitle]{enhanced,breakable,colback=fond,colframe=bleu,
  boxrule=0pt,leftrule=3pt,arc=0pt,left=8pt,right=8pt,top=6pt,bottom=6pt,
  title={\bfseries #1},coltitle=bleu,colbacktitle=fond,attach title to upper={\par\smallskip}}
\newtcolorbox{attention}[1][\wacwarningtitle]{enhanced,breakable,colback=fondalerte,colframe=alerte,
  boxrule=0pt,leftrule=3pt,arc=0pt,left=8pt,right=8pt,top=6pt,bottom=6pt,
  title={\bfseries #1},coltitle=alerte,colbacktitle=fondalerte,attach title to upper={\par\smallskip}}

% Code
\lstset{basicstyle=\ttfamily\small,columns=fullflexible,keepspaces=true,
  backgroundcolor=\color{fond},frame=l,framerule=3pt,rulecolor=\color{bleu},
  xleftmargin=9pt,framexleftmargin=6pt,aboveskip=0.8em,belowskip=0.8em,
  breaklines=true,breakatwhitespace=false,showstringspaces=false,
  literate={é}{{\'e}}1 {è}{{\`e}}1 {à}{{\`a}}1 {ê}{{\^e}}1 {ô}{{\^o}}1 {ç}{{\c c}}1
           {É}{{\'E}}1 {—}{{---}}1 {’}{{'}}1 {…}{{...}}1 {«}{{\guillemotleft}}1 {»}{{\guillemotright}}1}

% Breakable monospace. A path or a field name in \texttt is one word that no
% hyphenation pattern cuts: in a table cell it ran 190 pt past the column
% (ProgramData\Microsoft\NetFramework\BreadcrumbStore). As the \seqsplit of
% the EventsManager theme's tables, a line may break inside it, without a
% hyphen that would read as part of the name: after \ / . _ - : and between a
% lowercase letter or digit and a capital (AmcacheSha1|Covers|Bytes). Only
% inside \texttt — the classes are set once, the breaks switched on there.
\newXeTeXintercharclass\wacBreakAfter
\newXeTeXintercharclass\wacUpper
\newXeTeXintercharclass\wacLower
\XeTeXcharclass`\\=\wacBreakAfter
\XeTeXcharclass`\/=\wacBreakAfter
\XeTeXcharclass`\.=\wacBreakAfter
\XeTeXcharclass`\_=\wacBreakAfter
\XeTeXcharclass`\-=\wacBreakAfter
\XeTeXcharclass`\:=\wacBreakAfter
\count255=`A \loop\XeTeXcharclass\count255=\wacUpper\ifnum\count255<`Z\advance\count255 by 1\repeat
\count255=`a \loop\XeTeXcharclass\count255=\wacLower\ifnum\count255<`z\advance\count255 by 1\repeat
\count255=`0 \loop\XeTeXcharclass\count255=\wacLower\ifnum\count255<`9\advance\count255 by 1\repeat
\XeTeXinterchartoks\wacBreakAfter\wacUpper={\allowbreak}
\XeTeXinterchartoks\wacBreakAfter\wacLower={\allowbreak}
\XeTeXinterchartoks\wacBreakAfter 0={\allowbreak}
\XeTeXinterchartoks\wacLower\wacUpper={\allowbreak}
% \textbackslash is no character token: it carries its own break.
\newcommand{\wacbackslash}{\char92\allowbreak}
\let\wacplaintexttt\texttt
\renewcommand{\texttt}[1]{\wacplaintexttt{\XeTeXinterchartokenstate=1\let\textbackslash\wacbackslash#1}}

% Shortcuts
\newcommand{\wac}{\textsc{wac}}
\newcommand{\opt}[1]{\texttt{-{}-#1}}
\newcommand{\champ}[1]{\texttt{#1}}
\newcommand{\fichier}[1]{\texttt{#1}}

% ============================================================ COVER PAGE
% The cover images are 2:3 (1024 x 1536), taller than an A4 page once widened
% to it: centred, they lost 9 mm at the top — the CEC logo — and 9 mm at the
% bottom. Raised 4 mm above the top of the page instead: the logo, which starts
% 6 mm below the image's edge, stays whole, and the bottom row of text keeps
% 7 mm to the page's edge; only a decorative corner below it is cut.
% Usage: \coverpage{cover_FR.png}
\newcommand{\coverpage}[1]{%
  \begin{titlepage}
  \thispagestyle{empty}
  \begin{tikzpicture}[remember picture,overlay]
    \node[inner sep=0pt,anchor=north,yshift=4mm] at (current page.north)
      {\includegraphics[width=\paperwidth]{#1}};
  \end{tikzpicture}
  \end{titlepage}}
:
%   \usepackage[<language>]{babel}
%   \newcommand{\wacdoctitle}{...}      running header: "WAC — <title>"
%   \newcommand{\wacpageof}{...}        footer: "Page N <of> M"
%   \newcommand{\wacnotetitle}{...}     default title of the "note" box
%   \newcommand{\wacwarningtitle}{...}  default title of the "attention" box
% and AFTER it its \hypersetup (PDF title, author, subject).

% ---------------------------------------------------------------- fonts, language
\usepackage{fontspec}
\setmainfont{Noto Sans}
\setmonofont{DejaVu Sans Mono}[Scale=0.84]

% ---------------------------------------------------------------- layout
\usepackage[top=26mm,bottom=26mm,left=24mm,right=22mm,headheight=14pt]{geometry}
\usepackage{graphicx}
\graphicspath{{../}}
\usepackage{xcolor}
\usepackage{tikz}
\usepackage{booktabs,tabularx,longtable,array}
\usepackage{enumitem}
\usepackage{fancyhdr}
\usepackage{titlesec}
\usepackage[most]{tcolorbox}
\usepackage{listings}
\usepackage{url}
\usepackage{lastpage}
\usepackage[hidelinks,unicode,pdfpagelabels]{hyperref}

% Colours sampled from the cover page
\definecolor{marine}{HTML}{0A1A33}
\definecolor{bleu}{HTML}{2F6DB5}
\definecolor{accent}{HTML}{1F66D1}
\definecolor{fond}{HTML}{EEF3FA}
\definecolor{alerte}{HTML}{B3261E}
\definecolor{fondalerte}{HTML}{FBEDEC}
\definecolor{gris}{HTML}{5B6573}

\setlength{\parindent}{0pt}
\setlength{\parskip}{0.55em}
\linespread{1.08}
\setlist{itemsep=0.25em,topsep=0.35em}
\urlstyle{tt}
\renewcommand{\arraystretch}{1.25}
% Left-aligned table columns: justified narrow columns opened wide gaps
% between words.
\newcolumntype{L}[1]{>{\raggedright\arraybackslash}p{#1}}
% Monospaced file and field names cannot be hyphenated: let the line's spaces
% stretch instead.
\setlength{\emergencystretch}{3em}

% Headings
\titleformat{\chapter}[display]
  {\color{marine}\bfseries\raggedright}
  {\color{bleu}\large\MakeUppercase{\chaptertitlename}~\thechapter}
  {0.3em}
  {\Huge}[{\color{accent}\titlerule[2pt]}]
\titlespacing*{\chapter}{0pt}{0pt}{1.6em}
\titleformat{\section}{\color{marine}\Large\bfseries}{\color{bleu}\thesection}{0.8em}{}
\titleformat{\subsection}{\color{marine}\large\bfseries}{\color{bleu}\thesubsection}{0.7em}{}
\titleformat{\subsubsection}{\color{bleu}\normalsize\bfseries}{}{0em}{}

% Headers and footers
\pagestyle{fancy}
\fancyhf{}
\fancyhead[L]{\small\color{gris}WAC — \wacdoctitle}
\fancyhead[R]{\small\color{gris}\leftmark}
% Footer: "Page N of M" (the word set by \wacpageof), readable. A bare, small, light-grey number went
% unnoticed.
\newcommand{\piedpage}{\small\color{marine}\textbf{Page \thepage}\ \wacpageof\ \pageref*{LastPage}}
\fancyfoot[C]{\piedpage}
\renewcommand{\footrulewidth}{0.4pt}
\renewcommand{\footrule}{\hbox to\headwidth{\color{fond}\leaders\hrule height \footrulewidth\hfill}}
\renewcommand{\headrulewidth}{0.4pt}
\renewcommand{\headrule}{\hbox to\headwidth{\color{fond}\leaders\hrule height \headrulewidth\hfill}}
\renewcommand{\chaptermark}[1]{\markboth{#1}{}}
\fancypagestyle{plain}{\fancyhf{}\fancyfoot[C]{\piedpage}\renewcommand{\headrulewidth}{0pt}}

% Boxes
\newtcolorbox{note}[1][\wacnotetitle]{enhanced,breakable,colback=fond,colframe=bleu,
  boxrule=0pt,leftrule=3pt,arc=0pt,left=8pt,right=8pt,top=6pt,bottom=6pt,
  title={\bfseries #1},coltitle=bleu,colbacktitle=fond,attach title to upper={\par\smallskip}}
\newtcolorbox{attention}[1][\wacwarningtitle]{enhanced,breakable,colback=fondalerte,colframe=alerte,
  boxrule=0pt,leftrule=3pt,arc=0pt,left=8pt,right=8pt,top=6pt,bottom=6pt,
  title={\bfseries #1},coltitle=alerte,colbacktitle=fondalerte,attach title to upper={\par\smallskip}}

% Code
\lstset{basicstyle=\ttfamily\small,columns=fullflexible,keepspaces=true,
  backgroundcolor=\color{fond},frame=l,framerule=3pt,rulecolor=\color{bleu},
  xleftmargin=9pt,framexleftmargin=6pt,aboveskip=0.8em,belowskip=0.8em,
  breaklines=true,breakatwhitespace=false,showstringspaces=false,
  literate={é}{{\'e}}1 {è}{{\`e}}1 {à}{{\`a}}1 {ê}{{\^e}}1 {ô}{{\^o}}1 {ç}{{\c c}}1
           {É}{{\'E}}1 {—}{{---}}1 {’}{{'}}1 {…}{{...}}1 {«}{{\guillemotleft}}1 {»}{{\guillemotright}}1}

% Breakable monospace. A path or a field name in \texttt is one word that no
% hyphenation pattern cuts: in a table cell it ran 190 pt past the column
% (ProgramData\Microsoft\NetFramework\BreadcrumbStore). As the \seqsplit of
% the EventsManager theme's tables, a line may break inside it, without a
% hyphen that would read as part of the name: after \ / . _ - : and between a
% lowercase letter or digit and a capital (AmcacheSha1|Covers|Bytes). Only
% inside \texttt — the classes are set once, the breaks switched on there.
\newXeTeXintercharclass\wacBreakAfter
\newXeTeXintercharclass\wacUpper
\newXeTeXintercharclass\wacLower
\XeTeXcharclass`\\=\wacBreakAfter
\XeTeXcharclass`\/=\wacBreakAfter
\XeTeXcharclass`\.=\wacBreakAfter
\XeTeXcharclass`\_=\wacBreakAfter
\XeTeXcharclass`\-=\wacBreakAfter
\XeTeXcharclass`\:=\wacBreakAfter
\count255=`A \loop\XeTeXcharclass\count255=\wacUpper\ifnum\count255<`Z\advance\count255 by 1\repeat
\count255=`a \loop\XeTeXcharclass\count255=\wacLower\ifnum\count255<`z\advance\count255 by 1\repeat
\count255=`0 \loop\XeTeXcharclass\count255=\wacLower\ifnum\count255<`9\advance\count255 by 1\repeat
\XeTeXinterchartoks\wacBreakAfter\wacUpper={\allowbreak}
\XeTeXinterchartoks\wacBreakAfter\wacLower={\allowbreak}
\XeTeXinterchartoks\wacBreakAfter 0={\allowbreak}
\XeTeXinterchartoks\wacLower\wacUpper={\allowbreak}
% \textbackslash is no character token: it carries its own break.
\newcommand{\wacbackslash}{\char92\allowbreak}
\let\wacplaintexttt\texttt
\renewcommand{\texttt}[1]{\wacplaintexttt{\XeTeXinterchartokenstate=1\let\textbackslash\wacbackslash#1}}

% Shortcuts
\newcommand{\wac}{\textsc{wac}}
\newcommand{\opt}[1]{\texttt{-{}-#1}}
\newcommand{\champ}[1]{\texttt{#1}}
\newcommand{\fichier}[1]{\texttt{#1}}

% ============================================================ COVER PAGE
% The cover images are 2:3 (1024 x 1536), taller than an A4 page once widened
% to it: centred, they lost 9 mm at the top — the CEC logo — and 9 mm at the
% bottom. Raised 4 mm above the top of the page instead: the logo, which starts
% 6 mm below the image's edge, stays whole, and the bottom row of text keeps
% 7 mm to the page's edge; only a decorative corner below it is cut.
% Usage: \coverpage{cover_FR.png}
\newcommand{\coverpage}[1]{%
  \begin{titlepage}
  \thispagestyle{empty}
  \begin{tikzpicture}[remember picture,overlay]
    \node[inner sep=0pt,anchor=north,yshift=4mm] at (current page.north)
      {\includegraphics[width=\paperwidth]{#1}};
  \end{tikzpicture}
  \end{titlepage}}


\hypersetup{pdftitle={WAC — Dossier d'innocuité},
  pdfauthor={Centre d'Excellence Cyberdéfense Aérospatiale (CEC)},
  pdfsubject={Les traces que WAC laisse sur la machine examinée, et la procédure qui le vérifie}}

\begin{document}

% ============================================================ COVER PAGE
\pagenumbering{roman}
\coverpage{cover_harmlessness_FR.png}

% ============================================================ INFORMATION PAGE
\thispagestyle{empty}
\vspace*{\fill}
{\color{marine}\Large\bfseries WAC — Windows Artefacts Collector}\par
{\color{bleu}Dossier d'innocuité}\par
\medskip
\begin{tabular}{@{}ll@{}}
Version du logiciel & 2.0.0 \\
Édition du document & septembre 2026 \\
Éditeur & Centre d'Excellence Cyberdéfense Aérospatiale (CEC) \\
        & École de l'air et de l'espace, base aérienne 701, Salon-de-Provence \\
\end{tabular}

\bigskip
{\small\color{gris}Ce document établit ce que \wac{} fait à la machine qu'il
examine : chaque opération, la trace qu'elle laisse, et ce qu'il ne fait
jamais. Il donne les preuves de ces affirmations — la table des fonctions que
\fichier{WAC.exe} peut appeler, les contrôles qui refusent une compilation
fautive, la mesure faite sur une machine de test — et la procédure qui permet
à chacun de les refaire.}
\vspace*{2cm}
\clearpage

\pagenumbering{arabic}
\tableofcontents

% ============================================================ 1
\chapter{Objet}

\section{Ce que « innocuité » veut dire ici}

Une collecte se fait sur une machine qui est aussi un élément de preuve. Un
outil qui la modifie détruit une partie de ce qu'il vient chercher, et mêle ses
propres traces à celles de l'activité étudiée. \wac{} est dit \emph{inoffensif}
au sens suivant :

\begin{enumerate}
  \item il \textbf{n'écrit rien} sur les volumes de la machine examinée — ni
        fichier, ni registre, ni événement dans un journal ; ce qu'une
        surveillance en place (Sysmon, audit, EDR) enregistre de son exécution
        est décrit au chapitre~\ref{ch:traces} ;
  \item il \textbf{ne modifie aucune métadonnée} des fichiers qu'il lit : ni
        date d'accès, ni numéro de changement du journal NTFS ;
  \item il \textbf{n'émet rien sur le réseau} et ne résout aucun nom ;
  \item il \textbf{n'installe rien et ne persiste pas} : ni service, ni pilote,
        ni tâche, ni clé de démarrage ; à sa sortie il ne reste de lui que
        les traces que Windows produit pour tout programme lancé ;
  \item il \textbf{n'altère pas le fonctionnement} de la machine : aucun
        processus ouvert, aucun service arrêté ou reconfiguré, aucun privilège
        modifié ;
  \item ce qu'il ne peut pas éviter est \textbf{énuméré}, ici et, collecte par
        collecte, dans le journal
        \fichier{investigation.json}.
\end{enumerate}

\section{Ce que ce document ne couvre pas}

\begin{itemize}
  \item La justesse des données collectées : c'est l'objet du \emph{Dossier de
        validation}.
  \item L'exploitation courante de l'outil : c'est l'objet du \emph{Manuel
        d'utilisation}.
  \item La mise à jour des listes de confiance (\opt{update-trust}) : elle
        s'exécute sur le poste d'analyse, jamais sur la machine examinée, et
        c'est le seul mode qui ouvre une connexion réseau.
\end{itemize}

% ============================================================ 2
\chapter{Inventaire des traces}\label{ch:traces}

Les traces se rangent en trois familles : celles de l'intervention elle-même
(brancher un support, ouvrir une session), celles que Windows produit pour tout
programme lancé, et celles des opérations propres à \wac{}. Les deux premières
existent quel que soit l'outil ; elles sont listées pour que l'analyste les
reconnaisse au lieu de les prendre pour l'activité étudiée.

\section{L'intervention}

\begin{longtable}{@{}L{4.3cm}L{10.1cm}@{}}
\toprule
\bfseries Action & \bfseries Trace \\
\midrule
\endhead
Brancher le support de collecte & Clés \fichier{USBSTOR}, \fichier{USB},
  \fichier{MountedDevices}, \fichier{Windows Portable Devices},
  \fichier{VolumeInfoCache} ; fichier \fichier{setupapi.dev.log} ; journaux
  \fichier{Microsoft-Windows-Partition/Diagnostic} (1006),
  \fichier{Kernel-PnP}, \fichier{DeviceSetupManager}. \\
Ouvrir une session & Événements Security 4624 et 4672 (privilèges spéciaux) ;
  pour un compte jamais connecté sur la machine, création de son profil. \\
Élever l'invite de commandes & Exécution de \fichier{consent.exe} (contrôle de
  compte d'utilisateur) et ses propres traces d'exécution. \\
Lancer \fichier{cmd.exe} & Traces d'exécution de \fichier{cmd.exe}
  (section suivante). \\
\bottomrule
\end{longtable}

\section{L'exécution de WAC.exe par Windows}

Windows les produit pour tout programme ; \wac{} n'y peut rien et ne tente pas
de les effacer — les effacer serait écrire sur la machine.

\begin{longtable}{@{}L{4.3cm}L{10.1cm}@{}}
\toprule
\bfseries Artefact & \bfseries Trace \\
\midrule
\endhead
Prefetch & \fichier{C:\textbackslash Windows\textbackslash Prefetch\textbackslash
  WAC.EXE-<empreinte>.pf}, écrit par Windows quelques secondes après le
  lancement, puis à la fin --- pour un \fichier{WAC.exe} lancé du disque de la
  machine : lancé depuis une clé USB, mesuré sous Sysmon, il n'en a pas eu. \\
Amcache, ShimCache & Une entrée pour \fichier{WAC.exe} (chemin, taille, date,
  empreinte SHA-1 pour Amcache). Mesurée : l'entrée Amcache est écrite pendant
  la collecte, celle du ShimCache à l'arrêt de la machine, sous la forme
  \texttt{SIGN.MEDIA=<signature> WAC\textbackslash WAC.exe} pour une clé. \\
BAM & Une valeur pour \fichier{WAC.exe} sous la clé de l'utilisateur, datée de
  la fin de l'exécution --- là encore pour un lancement depuis le disque de la
  machine : aucune depuis la clé USB. \\
SRUM & Consommation de ressources du processus, agrégée par Windows. \\
Journaux & Security 4688 et 4689 si l'audit de création de processus est actif ;
  Sysmon~1 si Sysmon est installé ; télémétrie de l'antivirus ou de l'EDR
  (analyse de \fichier{WAC.exe}, et, pour un EDR, signalement possible de
  l'ouverture du volume en lecture brute). \\
Mémoire & Les pages de \wac{} occupent la mémoire pendant la collecte et
  peuvent en évincer d'autres vers le fichier d'échange. \\
\bottomrule
\end{longtable}

\begin{note}
\wac{} collecte ces mêmes artefacts : il se retrouve donc dans ses propres
résultats (Prefetch de la collecte suivante, BAM, Amcache). C'est attendu, et
daté : l'intervalle \champ{Collection.StartUtc} --- \champ{EndUtc} de
\fichier{investigation.json} permet de les isoler.
\end{note}

\section{Mesuré face à KAPE et aux outils EZ}

Sur un clone de la machine de test, figé avec Sysmon, chaque outil est lancé
depuis une clé USB à partir du même état ; le journal des modifications de C:,
borné avant et après l'outil, et Sysmon disent ce qu'il a écrit, et qui l'a
écrit (document de comparaison, chapitre « Traces laissées sur la machine
examinée »).

\begin{longtable}{@{}L{4.3cm}L{10.1cm}@{}}
\toprule
\bfseries Outil & \bfseries Ce que ses processus écrivent sur la machine \\
\midrule
\endhead
\wac{} & Rien : aucun fichier, aucune valeur de registre. \\
KAPE (cibles) & 4 valeurs \texttt{ZoneMap} d'Internet Settings dans la ruche
  de l'utilisateur. \\
KAPE et ses modules & Les mêmes valeurs, réécrites par 14 processus ; 28
  fichiers sur C: (miettes du runtime .NET, fichiers sentinelles dans le Temp
  de l'utilisateur). \\
Outils EZ en direct & Aucune valeur ni aucun fichier, mais des ouvertures en
  création de clés de stratégie et de réseau. \\
\bottomrule
\end{longtable}

Aucun n'a fait écrire d'événement dans les journaux de Windows eux-mêmes. Mais
tous sont enregistrés par la surveillance installée, comme le prévoit la
section précédente : Sysmon a noté, pour \wac{}, son processus, ses 2 lectures
brutes de volume et les 16\,382 fichiers qu'il a créés, tous sur la clé ;
l'audit de création de processus, inactif sur le clone, aurait inscrit un 4688
et un 4689. Pour tous, Windows ajoute une entrée Amcache pendant l'exécution et
une entrée ShimCache à l'arrêt ; tous lisent les volumes en brut. KAPE et les outils EZ
demandent en outre le runtime .NET~9 sur la machine examinée.

\section{Les opérations de WAC}

Chaque opération est consignée dans \fichier{investigation.json} avec, dans le
champ \champ{Footprint}, la trace qu'elle est censée laisser. Le tableau
ci-dessous les reprend toutes.

\begin{longtable}{@{}L{4.3cm}L{10.1cm}@{}}
\toprule
\bfseries Opération & \bfseries Trace \\
\midrule
\endhead
Lecture directe du volume (\path|\\.\C:|) & Un descripteur ouvert en
  \textbf{lecture seule}, sans cache (\texttt{FILE\_FLAG\_NO\_BUFFERING}),
  sur chaque volume effectivement lu. Aucun fichier n'est ouvert par le
  système : aucune date d'accès modifiée, aucune entrée du journal USN, aucun
  verrou. Si l'audit d'accès aux objets est configuré sur le volume :
  Security 4656/4663. Sans cache, la lecture ne remplace pas non plus le
  contenu du cache de fichiers de la machine. \\
Ruches du registre & Copiées par la lecture directe, puis lues sur la copie.
  Aucune clé du registre vivant n'est ouverte : \fichier{WAC.exe} n'importe
  aucune fonction \texttt{Reg*}. Le rejeu des journaux de transaction se fait
  sur la copie, avec un journal d'annulation. Dans la ruche SECURITY, sont lus
  les noms et dates des secrets LSA (jamais leur valeur), les comptes de
  connexions en cache (jamais leur contenu) et les droits des comptes. Dans la
  ruche SAM, les condensats des mots de passe sont déchiffrés et les questions
  de réinitialisation lues en clair : une collecte voulue, pour que WAE repère
  les mots de passe faibles. Tout se fait sur la copie, en lecture seule ; la
  machine examinée n'est pas touchée, et un condensat est à sens unique. \\
Journaux d'événements (\opt{events}) & Lus par la lecture directe : le service
  de journalisation n'est pas sollicité, et cette lecture ne produit aucun événement. \\
Volumes FAT et exFAT & Lus en brut par le périphérique du volume
  (\path|\\.\X:|), en lecture seule, volumes fixes comme amovibles : secteur
  de démarrage, FAT et répertoires, jamais un fichier ouvert, aucune date
  d'accès modifiée. Le volume qui reçoit la collecte n'est pas lu. \\
Binaires (\opt{binary}) & Lus par la lecture directe ; signatures et
  catalogues vérifiés en mémoire. Aucun magasin de certificats ouvert, aucune
  vérification de révocation en ligne. \\
Processus & Un instantané du noyau (\texttt{CreateToolhelp32Snapshot}) et une
  requête au service Terminal Services (\texttt{WTSEnumerateProcessesW})
  pour les propriétaires. Aucun processus n'est ouvert : \fichier{WAC.exe}
  n'importe pas \texttt{OpenProcess}. \\
Modules des processus & Un instantané par processus ; ceux que le noyau
  protège (LSASS, antivirus) le refusent, et le refus est consigné. \\
Sessions & Une requête en lecture à LSASS
  (\texttt{LsaEnumerateLogonSessions}, \texttt{LsaGetLogonSessionData}). \\
État des services & Une énumération en lecture seule du gestionnaire de
  services, droits \texttt{SC\_MANAGER\_ENUMERATE\_SERVICE |
  SC\_MANAGER\_CONNECT}. Aucun service n'est ouvert ; la configuration vient
  de la ruche copiée. \\
État réseau (\texttt{artefacts.network}) & Lecture en mémoire des tables de la
  pile TCP/IP par \fichier{iphlpapi.dll}, chargée à ce moment seulement. Aucun
  paquet émis, aucun nom résolu. \\
Horloge et démarrage & Requêtes en mémoire au noyau
  (\texttt{NtQuerySystemInformation}, \texttt{GetTickCount64}). \\
Réglages de la machine & Lecteur système, page de code, fuseau horaire, volumes
  et numéros de série : requêtes en mémoire, sans ouverture de fichier. \\
Jeton du processus & Lecture des privilèges de \wac{} lui-même, pour vérifier
  les droits d'administrateur. Aucun privilège n'est activé ni modifié. \\
Écriture des résultats & Sur le dossier \opt{output} uniquement, relatif au
  répertoire courant — le support de collecte. \\
\bottomrule
\end{longtable}

\begin{attention}
\wac{} écrit là où le désigne \opt{output}. Lancé depuis un dossier du disque
examiné, ou avec un \opt{output} qui y pointe, il y écrirait ses résultats :
fichiers créés, enregistrements du \path|$MFT| et du journal USN, blocs
écrasés. \textbf{Lancez-le depuis le support de collecte}, qui doit être le
répertoire courant.
\end{attention}

\section{Ce que WAC ne fait jamais}

\begin{itemize}
  \item Créer un cliché instantané (VSS), appeler COM ou WMI : aucune de ces
        bibliothèques n'est importée. Les clichés existants sont lus en brut
        sur le volume (catalogue et magasins de la zone de différences), par
        le même accès en lecture seule que le reste : aucun n'est monté,
        exposé ni demandé au service.
  \item Récupérer un fichier supprimé par un outil, ou par une écriture sur le
        volume : le contenu des fichiers supprimés est lu en brut
        (enregistrements de la \path|$MFT|, \path|$Bitmap|, clusters libres, blocs d'index des répertoires, et au niveau~3 tout l'espace libre ; volumes FAT et exFAT par leur périphérique),
        par le même accès en lecture seule que le reste, et copié sur le
        support de collecte seulement. Un cluster en usage n'est pas lu.
  \item Ouvrir une clé du registre vivant, ou le vider sur le disque
        (\texttt{RegFlushKey}).
  \item Charger un pilote, créer ou modifier un service, une tâche planifiée,
        une clé de démarrage.
  \item Ouvrir un processus, lire sa mémoire, injecter du code.
  \item Lancer un processus : aucune fonction qui en crée un n'est importée
        (\texttt{system}, \texttt{CreateProcess}, \texttt{ShellExecute},
        \texttt{WinExec}). L'effacement de l'écran au démarrage lançait
        \fichier{cmd.exe /c cls}, avec ses traces d'exécution ; la mesure
        comparée des traces laissées l'a montré, et l'écran est désormais
        effacé par l'API de la console.
  \item Activer un privilège (\texttt{AdjustTokenPrivileges}) : la lecture
        directe d'un volume ne demande que les droits d'administrateur.
  \item Ouvrir une connexion réseau, résoudre un nom de compte auprès du
        contrôleur de domaine (\texttt{LookupAccountSid}) : les comptes sont
        nommés hors ligne, depuis les ruches copiées.
  \item Écrire dans un fichier temporaire du système.
\end{itemize}

% ============================================================ 3
\chapter{Les preuves}\label{ch:preuves}

\section{La table des imports de WAC.exe}

Un programme Windows ne peut appeler une fonction du système que s'il
l'importe, ou s'il charge sa bibliothèque dynamiquement. La table des imports
de \fichier{WAC.exe} est donc la liste exhaustive de ce qu'il peut demander au
système — hors les deux chargements dynamiques décrits plus bas. Elle se lit
sans exécuter le programme, sous Linux :

\begin{lstlisting}
$ x86_64-w64-mingw32-objdump -p WAC.exe | grep -E 'DLL Name|^\s+[0-9a-f]+\s+[0-9]+ '
\end{lstlisting}

Hors \fichier{msvcrt.dll} (la bibliothèque C) :

\begin{description}[style=nextline,leftmargin=1.2em]
  \item[ADVAPI32.dll] \texttt{OpenSCManagerW}, \texttt{EnumServicesStatusExW},
    \texttt{CloseServiceHandle} (état des services) ; \texttt{OpenProcessToken},
    \texttt{GetTokenInformation}, \texttt{GetUserNameW}, \texttt{GetLengthSid},
    \texttt{IsValidSid} (le jeton de \wac{} lui-même). Aucune fonction
    \texttt{Reg*}, aucune fonction de modification de service.
  \item[Secur32.dll] \texttt{LsaEnumerateLogonSessions},
    \texttt{LsaGetLogonSessionData}, \texttt{LsaFreeReturnBuffer}.
  \item[WTSAPI32.dll] \texttt{WTSEnumerateProcessesW},
    \texttt{WTSFreeMemory}.
  \item[bcrypt.dll] \texttt{BCryptOpenAlgorithmProvider},
    \texttt{BCryptSetProperty}, \texttt{BCryptGetProperty}, \texttt{BCryptGenerateSymmetricKey},
    \texttt{BCryptDecrypt}, \texttt{BCryptDestroyKey},
    \texttt{BCryptCloseAlgorithmProvider} : le déchiffrement AES des condensats
    de la SAM, à la conversion seule. Du calcul en mémoire, sans fichier ni
    registre.
  \item[KERNEL32.dll] Fichiers, mémoire, fils d'exécution, heure, console,
    volumes, instantané des processus. Les fonctions qui écrivent ou
    suppriment (\texttt{DeleteFileW}, \texttt{MoveFileExW},
    \texttt{RemoveDirectoryW}, \texttt{CreateHardLinkW}, \texttt{SetEndOfFile})
    servent au dossier de sortie, ou appartiennent à la bibliothèque C++ liée
    statiquement (\texttt{std::filesystem}) sans que \wac{} les appelle
    ailleurs : chaque suppression ou renommage du code source porte sur le
    dossier de sortie, sa copie de travail ou son dossier d'attente.
\end{description}

\textbf{Chargements dynamiques.} Trois bibliothèques seulement, depuis
\path|System32| (\texttt{LOAD\_LIBRARY\_}\allowbreak\texttt{SEARCH\_SYSTEM32} ou,
sous un Windows~7 qui ne connaît pas ce drapeau, leur chemin absolu dans ce
dossier) :
\fichier{iphlpapi.dll}, pour l'état réseau, quand \texttt{artefacts.network}
est demandé ; \fichier{winhttp.dll}, pour \opt{update-trust} seul, et
\fichier{esent.dll}, qui rejoue les journaux d'une base ESE copiée (WebCache,
SRUM, BITS) dans la copie de travail, à la conversion seule — ni l'une ni
l'autre ne s'exécute sur la machine examinée : \wac{} y collecte, et n'y
convertit jamais.

\section{Les garde-fous de la compilation}

\fichier{build-windows.sh} refuse de produire un \fichier{WAC.exe} qui
enfreindrait ces règles. Après l'édition de liens, il inspecte la table des
imports et échoue :

\begin{itemize}
  \item si \fichier{WAC.exe} importe une fonction de résolution de noms qui
        interroge le domaine (\texttt{LookupAccountSid},
        \texttt{LookupAccountName}, \texttt{LsaLookupSids},
        \texttt{LsaLookupNames}) ;
  \item s'il importe statiquement une bibliothèque réseau
        (\fichier{winhttp}, \fichier{wininet}, \fichier{ws2\_32},
        \fichier{wsock32}, \fichier{iphlpapi}) ;
  \item s'il dépend d'une bibliothèque d'exécution MinGW : l'exécutable doit se
        suffire, pour ne rien déposer à côté de lui.
\end{itemize}

Avant même de compiler, il refuse un compilateur dont le \emph{modèle de
threads} n'est pas « win32 ». Un gcc en threads « posix » — celui de MSYS2
sous Windows — intègre winpthreads à \fichier{WAC.exe}, et avec lui 38
imports que \wac{} n'appelle jamais : \texttt{SetSystemTime},
\texttt{SetThreadContext}, \texttt{SuspendThread},
\texttt{TerminateProcess}, \texttt{OpenProcess},
\texttt{SetProcessAffinityMask}\ldots{} Ils ne changent rien à ce que fait
\wac{}, mais ils élargissent ce qu'il \emph{pourrait} demander au système, et
ce sont précisément les noms qu'un outil de défense ou un analyste remarque
sur la machine examinée. Mesuré : avec le gcc de Debian (Linux) ou celui de
MinGW-Builds (Windows), tous deux en threads « win32 » et sur msvcrt, ces
imports sont absents, et les deux tables ne diffèrent que de huit fonctions
d'entrée-sortie de la bibliothèque C. « win32 » nomme le modèle de threads,
pas une architecture : \fichier{WAC.exe} est x86\_64 dans tous les cas.

La compilation est reproductible : les mêmes sources donnent le même
\fichier{WAC.exe}, octet pour octet, et \fichier{SHA256SUMS} l'atteste.

\section{Le journal d'investigation}

Chaque collecte écrit, dans \fichier{investigation.json}, la liste ordonnée
de ses opérations, chacune avec sa cible, son résultat et son champ
\champ{Footprint}. Un écart entre ce que dit ce dossier et ce que dit le
journal d'une collecte réelle se voit donc collecte par collecte, et non
seulement lors d'une campagne de test.

% ============================================================ 4
\chapter{Procédure de validation}\label{ch:procedure}

Cette procédure est à refaire à chaque version, et à chaque modification qui
touche un appel au système. Les deux premières étapes se font sur le poste
d'analyse, sans machine Windows.

\section{Étape 1 : le binaire livré est celui des sources}

\begin{lstlisting}
$ ./build-windows.sh
$ sha256sum -c bin/SHA256SUMS
\end{lstlisting}

Attendu : la compilation se termine sans erreur, sans avertissement, et
affiche les trois contrôles : « Static C++ runtime », « No network name
resolution imported », « No network library imported ». \fichier{WAC.exe}
compilé et celui de \fichier{bin} ont la même empreinte.

\section{Étape 2 : la table des imports}

\begin{lstlisting}
$ x86_64-w64-mingw32-objdump -p bin/WAC.exe | grep -E 'DLL Name|^\s+[0-9a-f]+\s+[0-9]+ '
\end{lstlisting}

Attendu : les cinq bibliothèques et les fonctions de la
section~\ref{ch:preuves}, et aucune autre. Une fonction nouvelle
s'explique dans ce dossier avant la livraison.

\section{Étape 3 : la mesure sur une machine de test}

Le cycle automatique (\fichier{vmtest/run-wac-test.sh}, voir le \emph{Guide des
tests}) lance \wac{} dans une machine virtuelle Windows~11 et vérifie, sur des
fichiers témoins que \wac{} lit, qu'aucune de leurs métadonnées ne change :

\begin{enumerate}
  \item avant la collecte, le cycle crée des fichiers témoins que \wac{} lit
        forcément (un document à macro et un script, que \opt{binary}
        prélève) et relève, par PowerShell, leurs quatre dates et le numéro
        de leur dernier changement dans le journal NTFS
        (\texttt{fsutil usn readdata}) ;
  \item il lance la collecte ;
  \item il relève à nouveau les mêmes valeurs : elles doivent être identiques ;
  \item pour prouver que le contrôle sait échouer, il lit ensuite un témoin
        par l'API (\texttt{Get-Content}) et vérifie que son numéro USN, lui,
        change quand Windows met à jour sa date d'accès.
\end{enumerate}

Le résultat figure dans le compte rendu du cycle (\fichier{harmlessness.txt}).

\section{Étape 4 : le journal d'une collecte réelle}

Sur une collecte réelle, relire \fichier{investigation.json} : chaque
opération doit appartenir au tableau du chapitre~\ref{ch:traces}, avec le
\champ{Footprint} correspondant. Une opération inconnue de ce dossier est un
défaut à signaler.

% ============================================================ 5
\chapter{Limites}

\begin{description}[style=nextline,leftmargin=1.2em]
  \item[Dossier de sortie] \wac{} ne vérifie pas que \opt{output} désigne un
        autre volume que ceux qu'il examine : c'est à l'opérateur de le lancer
        depuis le support de collecte.
  \item[Surveillance] Un EDR ou une politique d'audit peut journaliser
        l'ouverture du volume en lecture brute, l'énumération des processus et
        des sessions. Ces journaux sont écrits par la machine, non par \wac{} ;
        ils sont attendus.
  \item[Mémoire] \wac{} ne capture pas la mémoire, mais il en occupe : une
        page évincée vers le fichier d'échange pendant la collecte n'est pas
        une trace sur le disque que \wac{} lit, mais elle change l'état de la
        mémoire que capturerait un outil lancé ensuite. Capturer la mémoire,
        si elle est nécessaire, \textbf{avant} \wac{}.
  \item[Mesure] L'étape~3 porte sur des fichiers témoins d'une machine
        virtuelle : elle établit le comportement de la lecture, pas l'absence
        de toute écriture de Windows lui-même pendant la collecte.
\end{description}

\end{document}
